% Source: https://artofproblemsolving.com/community/c3h2109018_weird_function_thing_i_thought_about
% Stable identity: aops-topic-2109018-q01
Let $g(y)=y(y-1)(y-2)$. The condition says $g\left(f(x)\right)=x$ for every $x$. Substituting $y=t+1$ makes $g$ symmetric:
\[g(t+1)=(t+1)\,t\,(t-1)=t^3-t .\]

\textbf{(a)} For each real $x$, the cubic equation $t^3-t-x=0$ has at least one real root, since every polynomial of odd degree with real coefficients does (it takes values of both signs and is continuous). So we may define
\[f(x)=1+\left(\text{the largest real root of }t^3-t=x\right).\]
Then $g\left(f(x)\right)=x$ for all $x$. For example, $f(0)=2$, $f(6)=3$ (since $2^3-2=6$) and $f(-6)=-1$ (the only real root of $t^3-t=-6$ is $t=-2$). Any choice of a real root for each $x$ gives such a function, so there are infinitely many.

\textbf{(b)} Suppose $f$ is continuous. If $f(x_1)=f(x_2)$, then $x_1=g\left(f(x_1)\right)=g\left(f(x_2)\right)=x_2$, so $f$ is one-to-one. The image $J=f(\mathbb R)$ of the line under a continuous function is an interval, and every real number $x=g\left(f(x)\right)$ lies in $g(J)$.

If $J$ were bounded above, say $J\subseteq(-\infty,M]$, then $g$ would be bounded above on $J$: indeed $g(y)\to-\infty$ as $y\to-\infty$, so $g(y)<0$ for $y$ below some $L$, and $g$ is bounded on $[L,M]$. That contradicts $g(J)=\mathbb R$. Likewise $J$ is not bounded below. So the interval $J$ is all of $\mathbb R$.

Now $0$, $1$ and $2$ are all in $J$, so $0=f(x_0)$ and $1=f(x_1)$ for some $x_0,x_1$. But then
\[x_0=g(0)=0=g(1)=x_1,\]
so $f(x_0)=f(x_1)$ with $0\ne1$, a contradiction. Hence $\boxed{\text{no such function is continuous}}$. (For example, the function in (a) jumps at $x=-\dfrac{2}{3\sqrt3}$, where the two larger roots of $t^3-t=x$ merge and disappear.)
